% Generated draft: requires mathematical and source review.
\documentclass[11pt]{article}
\usepackage[margin=25mm]{geometry}
\usepackage{fontspec}
\setmainfont{Latin Modern Roman}
\usepackage{amsmath,amssymb,mathtools,unicode-math}
\setmathfont{Latin Modern Math}
\usepackage{enumitem,needspace,xurl,hyperref}
\setlength{\parindent}{0pt}
\setlength{\parskip}{5pt}
\setlength{\emergencystretch}{3em}
\newcommand{\N}{\mathbb{N}}
\newcommand{\Z}{\mathbb{Z}}
\newcommand{\Q}{\mathbb{Q}}
\newcommand{\R}{\mathbb{R}}
\newcommand{\C}{\mathbb{C}}
\newcommand{\F}{\mathbb{F}}
\newcommand{\A}{\mathbb{A}}
\newcommand{\PP}{\mathbb{P}}
\newcommand{\E}{\mathbb{E}}
\newcommand{\eps}{\varepsilon}
\newcommand{\dd}{\,\mathrm{d}}
\begin{document}
\section*{UnsolvedMath ID: 11000262}
\textbf{Source problem code:} catalogue entry
\subsection{Definitions and mathematical statement}
\emph{Question 5.} Is there a homological definition of representations of the Birman--Wenzl--Murakami algebra? The source expresses belief that the answer is yes, but does not supply the intended homology spaces, coefficient/local-system data, maps, or factorization through the algebra. A precise construction must specify those data and verify the defining relations; no homological model is asserted here.

The accompanying source proposes an algebra \(Z_n\) generalizing Iwahori--Hecke and Birman--Wenzl--Murakami algebras. Retain \(R,q,t\) as its unspecified coefficients and parameters. It abbreviates \(\sigma_{i_1\cdots i_k}=\sigma_{i_1}\cdots\sigma_{i_k}\) and uses barred words for inverse notation. The extraction does not clearly distinguish inverse of the entire word from the product of inverses in the written order; bars below retain the source notation pending that clarification.

In the stated braid-group algebra \(R B_n\), it defines
\[
 X_2=q\bar\sigma_1+1-q-\sigma_1,
 \qquad
 X_k=\bigl(q^{k-1}\bar\sigma_{(k-1)\cdots1}
              -\sigma_{1\cdots(k-1)}\bigr)X_{k-1}\quad(3\leq k\leq n).
\]
Thus the displayed initial cases are
\[
 X_3=(q^2\bar\sigma_{21}-\sigma_{12})X_2,
 \qquad X_4=(q^3\bar\sigma_{321}-\sigma_{123})X_3.
\]
The proposed \(Z_n\) is the quotient by the displayed relations
\[
 (q\bar\sigma_2+1-q-\sigma_2)X_2
       =(q\bar\sigma_{21}-\sigma_{12})X_2,
\]
\[
 (q^2\bar\sigma_{32}-\sigma_{23})X_3
       =(q^2\bar\sigma_{321}-\sigma_{123})X_3,
\]
\[
 (q^3\bar\sigma_{432}-\sigma_{234})X_4
       =(q^3\bar\sigma_{4321}-\sigma_{1234})X_4,
\]
and the continuation terminating, as printed, with
\[
 (q^{n-1}\bar\sigma_{n\cdots2}-\sigma_{2\cdots n})X_n
       =(q^{n-1}\bar\sigma_{n\cdots1}-\sigma_{1\cdots n})X_n.
\]
There is an indexing conflict: standard \(B_n\) has generators only \(\sigma_1,\ldots,\sigma_{n-1}\), while the final relation uses \(\sigma_n\). The braid-index or relation-range convention must be repaired from the source before this is a complete algebra presentation; no silent replacement by \(B_{n+1}\) or truncation is made.

The source states that \(H_n(1,-q)\) is obtained by imposing \(X_2=0\), and the Birman--Wenzl--Murakami algebra by imposing \(X_3=0\) and \(\sigma_1X_2=tX_2\). These quotient identifications retain the unresolved parameter and index conventions. It calls the homological motivation unclear and unpublished at the time, suggests combinatorial computations, and introduces a following question about whether \(Z_n\) is bigger than the Birman--Wenzl--Murakami algebra; that following precise question is not supplied.
\subsection{Short English statement}
Seek a homological construction of Birman--Wenzl--Murakami representations. The source's proposed algebra \(Z_n\) and quotient formulas are retained as context, with unresolved inverse-word notation, coefficients, and a final braid-generator indexing conflict.
\subsection{Sources}
\begin{itemize}
\item Local UnsolvedMath registry, permanent ID 11000262, source code catalogue entry, Question 5, PDF page 322. The chapter header identifies S. Bigelow; no individual proposer field is supplied.
\item Farb, \emph{Problems on Mapping Class Groups and Related Topics}. Supplied URL: \url{https://www.math.uchicago.edu/~farb/papers/mcgbook.pdf}; recorded access 2026-07-29; extraction pdf-font-structure. The volume contains 329 explicitly typeset headings across chapters, with repeated labels disambiguated by page.
\item The source supplies the \(X_k\) recursion, relations, and claimed Hecke/Birman--Wenzl--Murakami quotients directly, without a separate citation. It explicitly describes its homological motivation as then unclear and unpublished.
\item Supplied status open and OPEN-TRIAGE. The summary says only ``Cannot verify.'' Current primary-literature status and redistribution terms for the publicly accessible author/publisher copy remain NEEDS\_REVIEW.
\end{itemize}
\end{document}
